\section{Introduction}\label{sec:introduction}

The Weisfeiler--Leman method compares graphs by repeatedly refining
colors of vertex tuples. Its dimension controls how much of a graph
each tuple can describe. A useful way to study this refinement is to
count graph homomorphisms from graphs of bounded treewidth: those
counts connect local information in a tree decomposition to the global
color distribution of the target graph.

We give an explicit reduction from compatible choices in finite domains
to Weisfeiler--Leman equivalence. Its main ingredients are two parity
lifts of one base graph and a procedure that extracts compatible choices
from any distinguishing homomorphism test with sufficiently small bags.
All output graphs are finite, simple, undirected and uncolored.
The construction is independent of the computational applications in
the last section.

\subsection{Refinement and simultaneous choices}

Fix an integer \(k\ge4\). Initially, an ordered tuple
\(\mathbf x\in V(X)^k\) has its \emph{atomic type} as color: the
equalities and adjacencies between all positions. Write
\(\mathbf x[i\leftarrow y]\) for replacement of coordinate \(i\) by \(y\).
In \emph{joint} refinement, also called the correlated replacement-vector
convention, a round records the old color \(C_r(\mathbf x)\) and
\[
 \{\!\{\,\bigl(C_r(\mathbf x[1\leftarrow y]),\ldots,
                  C_r(\mathbf x[k\leftarrow y])\bigr):
                       y\in V(X)\,\}\!\}.
\]
In \emph{separate} refinement it records the old color and the \(k\)
multisets
\(\{\!\{C_r(\mathbf x[i\leftarrow y]):y\in V(X)\}\!\}\), one for
each coordinate \(i\). Colors are these data, with common names across
graphs. Two graphs are \(k\)-WL equivalent when their tuple-color
histograms agree at every round. Because each round retains the old
color, it suffices to compare until the common color partition on the
two tuple sets stabilizes.
We also consider refinement of all tuples in a disjoint union
\(X_0\sqcup X_1\), comparing the histograms on its two pure tuple sets
\(V(X_0)^k\) and \(V(X_1)^k\).

The parameter used throughout the construction is
\begin{equation}\label{eq:bag-parameter}
 t=\begin{cases}
      k+1,&\text{joint refinement},\\
      k,&\text{separate refinement}.
    \end{cases}
\end{equation}

\begin{definition}[Simultaneous choices]\label{def:choice}
A \emph{choice system} consists of finite domains \(D_1,\ldots,D_t\)
and, for every \(i<j\), a finite label set \(L_{ij}\) and maps
\[
 \lambda_{ij}^i:D_i\longrightarrow L_{ij},
 \qquad
 \lambda_{ij}^j:D_j\longrightarrow L_{ij}.
\]
A \emph{successful choice} is a tuple
\((d_1,\ldots,d_t)\in\prod_iD_i\) satisfying
\(\lambda_{ij}^i(d_i)=\lambda_{ij}^j(d_j)\) for every \(i<j\).
Empty domains are allowed.
\end{definition}

For example, let \(D_i\) be a set of Boolean assignments on a finite
variable scope \(V_i\), and let both pair maps record restriction to
\(V_i\cap V_j\), with label set \(\{0,1\}^{V_i\cap V_j}\).
A successful choice is then a family of allowed partial assignments
that joins to one assignment on \(\bigcup_i V_i\).

Thus each pair of chosen elements must give the same answer to its
shared label test. The labels are input data for the construction;
they will not become vertex colors in the output. Write
\(X_0\equiv_k X_\star\) for equivalence under the chosen convention.

\begin{theorem}[Parity reduction]\label{thm:choice-reduction}
Fix \(k\ge4\), either refinement convention, and the corresponding
\(t\) in \eqref{eq:bag-parameter}. From every choice system on \(t\)
domains one can explicitly construct two uncolored simple graphs
\(X_0,X_\star\) of equal order such that
\[
 X_0\equiv_k X_\star
 \quad\Longleftrightarrow\quad
 \text{the choice system has no successful choice}.
\]
The assertion holds both with common color names on separate graphs
and with pure-tuple comparison in their disjoint union.

If every domain and label set has size at most \(A\), where \(A\ge1\),
their common order is at most \(A M_t\), where
\[
 M_t=t\,2^{\,t-2+3\binom{t-1}{2}}+12\binom t3.
\]
When a successful choice exists, the two graphs already have different
histograms after at most two joint rounds, or one separate round.
\end{theorem}

The order bound is a bound on the actual uncolored graphs; no promise
about distinguished fibers or preserved types is supplied to a decider.

\subsection{Context and the main ideas}

Weisfeiler and Leman introduced the refinement method in their study
of graph canonization~\citep{WL1968}. The parity examples of
Cai, F\"urer and Immerman~\citep{CFI1992} are pairs of nonisomorphic
vertex-colored graphs of maximum degree three that require WL dimension
linear in their order to distinguish. The present construction also compares
parity constraints, but its purpose is to make a successful selection
from prescribed finite domains exactly equivalent to a detectable
homomorphism-count difference.

Encoding local solutions of binary linear systems as graph vertices
also underlies the construction of Atserias, Man\v{c}inska, Roberson,
\v{S}\'amal, Severini and Varvitsiotis
\citep[author version, Section~6.2]{AMRSSV2019}.
Their graphs join assignments that disagree on shared variables and
compare a system with the system obtained by setting all right-hand
sides to zero. Roberson's parity graphs instead use agreement along
adjacent constraint types.

The lift-counting argument has a close predecessor in Roberson's parity
graphs and oddomorphisms~\citep[preprint version, Definition~3.1,
Theorem~3.6 and Lemma~3.8]{Roberson2026}. We use a typed-template adaptation of that
construction: a lift retains both a base-graph vertex and the parity
tag of its template type. The comparison of homogeneous and affine
solution sets, and its dual parity certificate, follow the same
framework. We prove these facts locally in the form needed for arbitrary
projections into our base graph. The additional work is to encode pair
labels by helpers and extract a successful choice from every
bounded-bag obstruction.

Dvo\v{r}\'ak~\citep{Dvorak2010} and Dell, Grohe and
Rattan~\citep{DGR2018} connect Weisfeiler--Leman refinement with
homomorphism counts from graphs of bounded treewidth. We use their
rooted-count viewpoint and give the exact two interfaces needed here:
equality of all counts with bags of size at most \(t\) implies
\(k\)-WL equivalence, and one explicit template is detected in the
number of rounds stated in Theorem~\ref{thm:choice-reduction}.
The proof keeps the difference between \(t=k+1\) and \(t=k\) visible,
including when refinement is run on the whole disjoint union.

A graph homomorphism is a vertex map preserving adjacency; write
\(\hom(F,X)\) for the number of maps from \(F\) to \(X\).
Our template has \(t\) main vertices forming a clique. For each pair
\(i<j\) and each third main vertex \(\ell\), it has a helper adjacent
to \(i,j,\ell\). The base graph replaces main type \(i\) by its domain
\(D_i\) and this helper type by a copy of \(L_{ij}\). Main edges and the
two helper edges toward \(i,j\) enforce label agreement; the helper
edge toward \(\ell\) is unrestricted.

The parity lifts replace each base vertex by bit vectors indexed by
its neighboring template types. Each vector has a prescribed parity,
and two vectors over a base edge must agree in the coordinates facing
each other. The two lifts differ only in the parity at main type \(1\).
For a fixed projected homomorphism into the base, the lifts to \(X_0\)
form the kernel of a binary linear system. Those to \(X_\star\) form
either a translate of the kernel or the empty set. Every projection
therefore contributes a nonnegative difference to
\(\hom(F,X_0)-\hom(F,X_\star)\). A single positive contribution cannot
be cancelled by other, possibly non-type-preserving homomorphisms.

A successful choice supplies such a contribution from the template
itself. Conversely, a distinguishing test gives a dual parity
certificate consisting of weights on its vertices and edges. For a
test with bags of size at most \(t\), these weights orient each edge of
a tree decomposition toward an odd side. A sink bag contains one
representative of every main type. This odd-side orientation is closely
related to Neuen and Seppelt's separator argument for oddomorphisms
\citep[Lemma~4.4 and proof of Theorem~4.1]{NS2026}.
The remaining issue is the representatives'
compatibility: a bag need not induce a clique, and its representatives
need not have nonzero weight. The helper associated with a mismatched
pair transfers its weighted discrepancy to whichever main type is
absent from a separator. The resulting boundary identity rules out
the mismatch, even with repeated projected images and zero weights.

Section~\ref{sec:wl-tests} proves the two homomorphism interfaces.
Section~\ref{sec:parity-lifts} constructs the graphs and establishes
noncancellation and the explicit positive witness.
Section~\ref{sec:separator} proves extraction and completes
Theorem~\ref{thm:choice-reduction}.
Section~\ref{sec:complexity} applies the construction to satisfiability.

\subsection{Conditional time bounds and their scope}

For fixed \(k\), direct refinement decides equivalence in
\(n^{O(k)}\) time. The complexity question is whether this growing
exponent can be avoided by a different algorithm.

The \emph{positive exponential-rate ETH} of Impagliazzo and
Paturi~\citep{IP2001} asserts that some \(\delta>0\) rules out a
deterministic \(3\)-SAT algorithm with running time
\(2^{\delta N}\operatorname{poly}(L)\), where \(N\) is the number of
variables and \(L\) the formula encoding length. This is the hypothesis
for the following fixed-dimension application.

\begin{theorem}[Fixed-dimension ETH consequence]\label{thm:eth}
Assume positive exponential-rate ETH. There are constants \(c>0\)
and \(K\ge4\) such that, for every fixed integer \(k\ge K\), no
deterministic algorithm decides \(k\)-WL equivalence of two
\(n\)-vertex graphs in time \(O(n^{ck})\).
This holds for both joint and separate refinement. The inputs are
explicit adjacency matrices of uncolored simple undirected graphs,
and time is measured on a multitape Turing machine.
\end{theorem}

The hypothetical algorithm and its multiplicative constant may depend
arbitrarily on the fixed \(k\); the constants \(c,K\) do not.
Theorem~\ref{thm:eth} therefore excludes even an ineffectively selected
family of solvers whose running-time exponents are \(o(k)\).
Under the weaker premise excluding one
uniform \(2^{o(N)}\operatorname{poly}(L)\)-time \(3\)-SAT algorithm, we
instead prove the uniform \(f(k)n^{o(k)}\) exclusion in
Proposition~\ref{prop:uniform}. Its algorithm receives \(k\) as input.
These hypotheses are kept separate throughout the parameter analysis.

Bounds on the number of refinement rounds concern a different question.
For each fixed \(k\ge2\), Grohe, Lichter, Neuen and Schweitzer
\citep[preprint version, Theorem~1 and Corollary~34]{GLNS2025}
construct colored graph pairs requiring \(\Omega(n^{k/2})\) rounds of
joint refinement. Such a bound constrains the refinement process itself;
it does not rule out a faster algorithm deciding its stable equivalence
relation.
The time exclusions proved here already apply to comparing tuple-color
histograms after two joint rounds or one separate round, under the same
respective hypotheses; Section~\ref{sec:complexity} gives the deduction.

Lichter, Ra{\ss}mann and Schweitzer~\citep[Section~6]{LRS2025}
expressed the expectation that neither equivalence nor identification
admits \(n^{o(k)}\) time and suggested lower bounds based on ETH or
its strong form.
The application above concerns the equivalence question.
The companion article \emph{Unconditional time lower bounds for
Weisfeiler--Leman equivalence}~\citep[Theorem~1.1]{CompanionUnconditional}
proves a stronger time bound: for every sufficiently large fixed \(k\)
and every correct deterministic decider, an absolute exponent linear
in \(k\) is required at every sufficiently large graph order.
It also obtains connected graphs of diameter at most two.
Its unconditional theorem supersedes the conditional lower bound for
stable equivalence in Theorem~\ref{thm:eth}.
The contribution of this article is the explicit choice reduction
and its helper-based bounded-bag extraction, using the parity-lift and
noncancellation framework described above. No part of its proof uses
that companion theorem.

For the ETH application, the sparsification lemma of Impagliazzo,
Paturi and Zane~\citep{IPZ2001} reduces a formula to sparse disjuncts.
Partitioning their clauses into \(t\) groups gives domains of size
\(2^{O(N/t)}\): each domain consists of the assignments satisfying
one group, and pair labels are restrictions to common variables.
We retain the sparsification rate, its dependent density constant,
the explicit matrix-writing cost and the simulator's input cost in
the full calculation. Extensions to other machine models require an
absolute-power Turing simulation, as specified in
Section~\ref{sec:complexity}.
